% speaker_notes.tex -- separate speech notes for the Beamer defence deck.
% Compile separately from main.tex to produce speaker_notes.pdf.

\documentclass[11pt,a4paper]{article}

\usepackage[T1]{fontenc}
\usepackage{lmodern}
\usepackage[margin=22mm]{geometry}
\usepackage{xcolor}
\usepackage{enumitem}
\usepackage{amsmath}
\usepackage{amssymb}
\usepackage[hidelinks]{hyperref}

\definecolor{AcademicNavy}{HTML}{17324D}
\definecolor{AcademicMuted}{HTML}{5F6B76}
\definecolor{AcademicRule}{HTML}{CBD5E1}

\setlength{\parindent}{0pt}
\setlength{\parskip}{4pt}
\emergencystretch=2em
\setlist[itemize]{leftmargin=*,topsep=2pt,itemsep=2pt}
\setlist[enumerate]{leftmargin=*,topsep=2pt,itemsep=2pt}

\newcounter{slideno}
\newenvironment{slidenote}[1]{%
  \refstepcounter{slideno}%
  \section*{\color{AcademicNavy}Slide \theslideno. #1}%
  \addcontentsline{toc}{section}{Slide \theslideno. #1}%
}{%
  \vspace{0.6em}%
  {\color{AcademicRule}\hrule}%
  \vspace{0.8em}%
}

\newcommand{\Purpose}{\vspace{0.2em}\textbf{Purpose:}\par}
\newcommand{\WhatToSay}{\vspace{0.2em}\textbf{What to say:}\par}
\newcommand{\KeyPoints}{\vspace{0.2em}\textbf{Key points:}\par}
\newcommand{\Transition}{\vspace{0.2em}\textbf{Transition to next slide:}\par}

\title{\color{AcademicNavy}Speaker Notes\\
AI-Guided Prediction of Thermoelastic Wave Propagation in Geological Media}
\author{Diploma defence speech notes}
\date{Bachelor's diploma defence, 2026}

\begin{document}

\maketitle
\tableofcontents
\newpage

\begin{slidenote}{AI-Guided Prediction of Thermoelastic Wave Propagation in Geological Media}
\Purpose
Open the defence, introduce the team, and set the research-prototype scope.

\WhatToSay
Good morning. The topic of our diploma work is \textit{AI-Guided Prediction of Thermoelastic Wave Propagation in Geological Media}. The work was prepared by Nurshanov Dias, Askarov Insar, Kassymkhanov Timur, and Kaidrakhmetova Dinara under the supervision of Dr. Bakhyt N. Alipova. The main idea is to study how artificial intelligence can support directional prediction of thermoelastic wave behaviour in simplified geological media. From the start, I would like to define the scope clearly: this is a research prototype and a comparative framework. It is not a field-calibrated geophysical simulator and it is not proposed as a replacement for full numerical modelling.

\KeyPoints
\begin{itemize}
  \item The defence concerns AI-assisted prediction of thermoelastic wave propagation.
  \item The task is source--probe directional prediction in simplified geological media.
  \item The result is a comparative research framework, not an industrial simulator.
\end{itemize}

\Transition
I will begin with the motivation and aim of the work.
\end{slidenote}

\begin{slidenote}{Motivation and Aim}
\Purpose
Explain why the topic matters and state the aim without overclaiming the prototype.

\WhatToSay
Temperature changes in rocks can affect stress, deformation, and wave propagation. This is relevant for geothermal systems, underground engineering, and geophysical interpretation, where thermal and mechanical effects can interact. The aim of the thesis is to develop and justify an AI-guided framework for directional prediction of thermoelastic wave propagation in geological media under simplified effective-medium assumptions. The important word here is ``framework''. We are not trying to model a real field site with all geological complexity. Instead, we create a controlled environment where different AI surrogate models can be compared under the same source--probe prediction task.

\KeyPoints
\begin{itemize}
  \item Thermal changes can create mechanical effects in rock.
  \item The aim is a directional prediction framework under simplified assumptions.
  \item The scope is a research prototype, not a field-validated simulator.
\end{itemize}

\Transition
Next, I will show the geological media and effective parameters used in the model.
\end{slidenote}

\begin{slidenote}{Geological Media We Model}
\Purpose
Show how geological materials enter the prototype through effective scalar parameters.

\WhatToSay
This slide shows representative effective parameters for four geological media: granite, basalt, sandstone, and limestone. The table includes density, compressional wave velocity, thermal conductivity, heat capacity, and thermal expansion where available. These values are used as effective scalar parameters under a locally homogeneous isotropic approximation. This means the model does not explicitly resolve pores, cracks, bedding, fluids, or anisotropy. The diagram on the right shows the simplified two-dimensional plane-strain domain with a thermal source point and a probe point. The line between them defines the direction of prediction.

\KeyPoints
\begin{itemize}
  \item Rock types differ in mechanical and thermal properties.
  \item The prototype uses effective isotropic material parameters.
  \item The source--probe geometry defines the directional query.
\end{itemize}

\Transition
After defining the materials, I will introduce the coupled thermoelastic system.
\end{slidenote}

\begin{slidenote}{Coupled Thermoelastic System}
\Purpose
Present the theoretical backbone of the work in a compact, defensible way.

\WhatToSay
The coupled thermoelastic system combines three ideas. First, the constitutive law relates stress to elastic strain and temperature perturbation. The term $\gamma=3K\alpha$ represents thermoelastic coupling, where $K$ is the bulk modulus and $\alpha$ is thermal expansion. Second, the equation of motion links stress gradients to acceleration and displacement. Third, the heat equation describes the evolution of temperature and includes coupling with volumetric deformation. In physical terms, heat changes temperature, temperature produces thermal strain, strain contributes to stress, and stress drives displacement and wave-like response.

\KeyPoints
\begin{itemize}
  \item Stress contains elastic and thermal contributions.
  \item Motion is driven by stress gradients.
  \item Heat transfer and deformation are coupled through the thermoelastic term.
\end{itemize}

\Transition
Now I will translate this physical model into the two-dimensional directional prediction task.
\end{slidenote}

\begin{slidenote}{2D Directional Prediction Task}
\Purpose
Define the practical prediction query and make the mathematical scope explicit.

\WhatToSay
The practical task is two-dimensional directional prediction. The model domain is $\Omega\subset\mathbb{R}^{2}$, and we assume plane strain, small deformation, and linear isotropic thermoelasticity. The primary fields are temperature $T(x,t)$ and in-plane displacement $u(x,t)=(u,v)$. A source point $x_s$ and a probe point $x_p$ define the vector $d=x_p-x_s$, the distance $L$, the unit direction $\hat d$, and the azimuth angle $\phi$. The model mapping $\mathcal{F}:X\to\widehat{Y}$ means that material, geometry, and time inputs are mapped to a predicted response. The scope is directional response, not complete field-scale geological simulation.

\KeyPoints
\begin{itemize}
  \item The task is defined by a source point and a probe point.
  \item Distance, unit direction, and azimuth are explicit geometric quantities.
  \item The prediction maps material, geometry, and time inputs to selected outputs.
\end{itemize}

\Transition
To train and compare the models, the work uses a controlled synthetic dataset.
\end{slidenote}

\begin{slidenote}{COMSOL-Based Synthetic Dataset}
\Purpose
Explain where the reference data come from and why they are synthetic.

\WhatToSay
The synthetic dataset was generated in COMSOL Multiphysics using a coupled thermoelastic model. The setup is a two-dimensional plane-strain geological domain with a size of $1\,\mathrm{m}\times1\,\mathrm{m}$. The initial temperature is $293.15\,\mathrm{K}$ and the source temperature is $1500\,\mathrm{K}$. From this model we export temperature, displacement, stress, strain, and velocity fields. These are not field measurements. They are controlled numerical reference data, which makes them suitable for comparing several surrogate models under the same assumptions.

\KeyPoints
\begin{itemize}
  \item The dataset is COMSOL-derived and synthetic.
  \item The domain is a simplified 2D plane-strain setting.
  \item Exported fields provide reference data for surrogate comparison.
\end{itemize}

\Transition
The next slide shows the mesh and numerical domain used for these exports.
\end{slidenote}

\begin{slidenote}{COMSOL 2D Model and Mesh}
\Purpose
Show the numerical domain and explain why mesh structure matters.

\WhatToSay
This slide shows the two-dimensional COMSOL model and finite-element mesh. The domain represents a simplified section of geological medium with an internal heated source. The mesh is refined near the source because this is where the strongest temperature, displacement, and stress gradients are expected. This mesh is important because it defines the spatial nodes and local structure from which the COMSOL fields are exported. Later, this same structure is useful for mesh-aware modelling, especially for the MeshGraphNet route.

\KeyPoints
\begin{itemize}
  \item The simulation is a simplified two-dimensional domain.
  \item Mesh refinement is concentrated around the thermal source.
  \item The mesh provides the spatial structure for exported fields.
\end{itemize}

\Transition
Now I will show the main simulated thermoelastic fields.
\end{slidenote}

\begin{slidenote}{COMSOL Simulation Results}
\Purpose
Connect the synthetic thermal source with temperature, displacement, and stress outputs.

\WhatToSay
The three images show the main result fields from the COMSOL setup. The first image is the temperature field: local heating creates a high-temperature region near the source while the surrounding rock remains closer to the initial temperature. The second image shows total displacement caused by thermal expansion. The third image shows von Mises stress, with stress concentration near the heated region. Together, these fields demonstrate the coupled thermal and mechanical response used as the reference basis for the neural surrogate models.

\KeyPoints
\begin{itemize}
  \item Local heating produces a localized temperature disturbance.
  \item Thermal expansion creates displacement.
  \item Stress concentration confirms coupled thermoelastic behaviour in the synthetic setup.
\end{itemize}

\Transition
After simulation, the exported fields must be represented in a model-ready format.
\end{slidenote}

\begin{slidenote}{Data Representation and Preprocessing}
\Purpose
Explain how COMSOL exports are converted into inputs for different model families.

\WhatToSay
The COMSOL simulation results are exported as CSV files for each field and time step. These exports are aligned into a common node-based representation. Dynamic fields are assembled into a spatio-temporal tensor $\mathbf{S}\in\mathbb{R}^{N_t\times N_p\times F_d}$, where $N_t$ is the number of time steps, $N_p$ is the number of spatial points, and $F_d$ is the number of dynamic field channels. Static material and geometry features are stored separately. Before training, each feature channel is normalized with the usual mean and standard deviation formula. This gives a common processed dataset for coordinate-based, grid-based, graph-based, and attention-based routes.

\KeyPoints
\begin{itemize}
  \item CSV exports are aligned by nodes and time steps.
  \item Dynamic fields form a spatio-temporal tensor.
  \item Normalization makes model training and comparison more stable.
\end{itemize}

\Transition
Next, I will describe the comparative prototype framework.
\end{slidenote}

\begin{slidenote}{Comparative Prototype Framework}
\Purpose
Explain that all model routes are evaluated under one shared workflow.

\WhatToSay
The prototype uses one workflow for all model routes. The frontend, backend orchestrator, material catalog, and model services follow one shared input--output contract. This is important for fair comparison: each route receives the same source--probe geometry, the same material presets, and the same normalized response format. The goal is not to present a production platform. The goal is to make a consistent research environment where different AI-guided prediction methods can be compared under the same conditions.

\KeyPoints
\begin{itemize}
  \item All routes use comparable inputs.
  \item Outputs are normalized into the same response structure.
  \item The framework supports fair model comparison.
\end{itemize}

\Transition
Before discussing individual models, I will define the shared comparison contract.
\end{slidenote}

\begin{slidenote}{Shared Model Comparison Contract}
\Purpose
Clarify the inputs, outputs, and interpretation shared by all surrogate routes.

\WhatToSay
All routes are treated as surrogate modelling approaches for the same AI-guided thermoelastic prediction task. The input side includes material parameters, source--probe geometry, observation time, and processed field features. The output side includes estimated temperature, displacement components, travel time, and diagnostic response quantities. The comparison basis is the same material presets, geometry contract, response format, and analytical travel-time sanity check. This keeps the comparison fair and avoids presenting the models as complete numerical solvers.

\KeyPoints
\begin{itemize}
  \item The comparison uses one shared prediction contract.
  \item Inputs and outputs are kept consistent across models.
  \item The routes are surrogate models, not equivalent physical solvers.
\end{itemize}

\Transition
I will start the model discussion with the physics-informed baseline.
\end{slidenote}

\begin{slidenote}{PINN: Physics-Informed Baseline}
\Purpose
Introduce the PINN route and explain why it is the main physical baseline.

\WhatToSay
PINN stands for Physics-Informed Neural Network. In this framework, it is the explicitly physics-informed route because its training objective can include thermoelastic residuals in addition to supervised data terms. It predicts temperature and displacement fields and serves as the physics-informed baseline for comparison with more data-driven routes. The interpretation must still be balanced: PINN is a neural surrogate whose loss encourages physical consistency. It is not an exact numerical solver and it is still limited by the simplified two-dimensional source--probe setup.

\KeyPoints
\begin{itemize}
  \item PINN combines supervised data with physical residual terms.
  \item It predicts temperature and displacement fields.
  \item It is the strongest physics-informed baseline in the comparison.
\end{itemize}

\Transition
The next slide shows the PINN input contract and architecture.
\end{slidenote}

\begin{slidenote}{PINN Input and Architecture}
\Purpose
Explain what the PINN receives, what it predicts, and why the architecture is differentiable.

\WhatToSay
The public neural contract maps coordinates, time, and material parameters to temperature and displacement outputs. The input includes coordinates, time, Young's modulus, Poisson's ratio, density, thermal expansion, conductivity, and heat capacity. The output includes temperature and displacement components. The architecture can be viewed as an MLP-based route, with a more developed version using separate coordinate and material encoders, a fused residual trunk, and separate heads for temperature and displacement. Smooth, differentiable outputs are important because physics residuals require spatial and temporal derivatives through automatic differentiation.

\KeyPoints
\begin{itemize}
  \item Inputs are geometry, time, and material properties.
  \item Outputs are temperature and displacement components.
  \item Differentiability allows residuals to be evaluated from derivatives.
\end{itemize}

\Transition
With the architecture defined, I will explain the composite loss.
\end{slidenote}

\begin{slidenote}{PINN Composite Loss and Residuals}
\Purpose
Show how data fitting and physical consistency are combined in the PINN route.

\WhatToSay
The PINN loss combines supervised field loss, velocity consistency, elastic wave residuals, and coupled thermal residuals. The supervised term compares predicted fields against COMSOL-derived reference data. The velocity term supports dynamic interpretation from time derivatives. The wave and thermal residuals encourage consistency with the thermoelastic equations introduced earlier. This is the main reason PINN is treated differently from purely supervised baselines. At the same time, the slide also states a limitation: dedicated boundary-condition and initial-condition losses are not yet enforced, so the result must still be validated carefully.

\KeyPoints
\begin{itemize}
  \item The loss combines supervised, velocity, elastic, and thermal terms.
  \item Residuals are computed using automatic differentiation.
  \item Boundary and initial losses are a current limitation.
\end{itemize}

\Transition
The second model family is the Fourier Neural Operator.
\end{slidenote}

\begin{slidenote}{Fourier Neural Operator}
\Purpose
Introduce FNO as the structured-grid neural operator route.

\WhatToSay
The Fourier Neural Operator is included as a field-to-field neural operator. Instead of predicting only a scalar output, it learns a mapping from an input field space to an output field space. In the practical two-dimensional setup, the input tensor has batch, channel, and two spatial dimensions, and the output tensor has the same spatial grid with selected output channels. For the three-channel thermoelastic setup, the outputs are temperature and two in-plane displacement components. FNO is useful because it can exploit regular grid structure, but in the current prototype it remains a supervised neural-operator branch rather than an explicitly physics-informed solver.

\KeyPoints
\begin{itemize}
  \item FNO learns a mapping between spatial fields.
  \item The practical setup uses two-dimensional tensors.
  \item The output channels represent temperature and in-plane displacement.
\end{itemize}

\Transition
The next slide shows the main FNO processing pipeline.
\end{slidenote}

\begin{slidenote}{FNO Pipeline}
\Purpose
Summarize the processing stages of the FNO route.

\WhatToSay
The FNO pipeline maps an input field to an output field through several stages. First, a lifting layer maps input channels into latent features. Then spectral blocks process the latent field. Finally, a projection layer maps the latent representation back to physical output channels. Inside the spectral block, the model uses Fourier transform, retained modes with learned weights, inverse Fourier transform, and a local convolution branch. In simple terms, the spectral part helps capture broader field patterns, while the local branch preserves nearby spatial corrections.

\KeyPoints
\begin{itemize}
  \item FNO uses lifting, spectral processing, and projection.
  \item Fourier operations capture global spatial structure.
  \item Local convolution helps represent short-range details.
\end{itemize}

\Transition
Now I will explain the spectral block more directly.
\end{slidenote}

\begin{slidenote}{FNO Spectral Block}
\Purpose
Explain the Fourier-space operation and retained modes at an intuitive level.

\WhatToSay
The spectral block updates the latent field by combining a spectral branch and a local convolution branch. The spectral branch transforms the field into Fourier space, applies a learned spectral weight tensor to selected coefficients, and returns the result to physical space using the inverse transform. Only a limited number of low-frequency modes are retained. This reduces computational cost and controls model capacity. The important interpretation is that FNO learns global and local field structure from data, but these operations do not directly enforce the thermoelastic PDE residuals.

\KeyPoints
\begin{itemize}
  \item The spectral branch works in Fourier space.
  \item Retained modes control spectral resolution and cost.
  \item The block is data-driven, not a PDE residual calculation.
\end{itemize}

\Transition
The next slide gives the training objective and the current caveat for FNO.
\end{slidenote}

\begin{slidenote}{Training and Evaluation of FNO}
\Purpose
State how FNO is trained and make the current calibration limitation transparent.

\WhatToSay
The FNO route is trained with a supervised objective combining mean squared error and relative error. Normalization is computed on the training set and reused during validation, testing, and inference. Evaluation includes visual field comparison, channel-wise errors, relative $L^2$ error, and physical sanity checks. The most important caveat is that the current FNO prototype still requires output-scale recalibration before it can be treated as a physically reliable predictor. Therefore, in the results section it should be interpreted as a diagnostic branch rather than as a final calibrated model.

\KeyPoints
\begin{itemize}
  \item FNO training uses supervised MSE and relative error terms.
  \item Channel normalization is reused consistently.
  \item The current output scale needs recalibration.
\end{itemize}

\Transition
The third model family is MeshGraphNet.
\end{slidenote}

\begin{slidenote}{MeshGraphNet: Graph Formulation}
\Purpose
Introduce MeshGraphNet as the mesh-aware surrogate route.

\WhatToSay
MeshGraphNet represents the COMSOL computational mesh as a graph. The mesh nodes form the set $\mathcal{V}$ and neighbouring connections form the edge set $\mathcal{E}$. Each node contains coordinates, current physical state, material parameters, and scenario conditions. Each edge contains relative geometry, such as the vector and distance between neighbouring nodes. This is useful because COMSOL data are naturally mesh-based. MeshGraphNet can preserve nonuniform mesh connectivity more naturally than a regular grid model, although it is still a supervised surrogate rather than a full physical solver.

\KeyPoints
\begin{itemize}
  \item The mesh is represented as nodes and edges.
  \item Node features carry state, material, and scenario information.
  \item Edge features carry local geometric relations.
\end{itemize}

\Transition
The next slide shows the input and output structure in more detail.
\end{slidenote}

\begin{slidenote}{MeshGraphNet: Inputs and Outputs}
\Purpose
Clarify what is stored at nodes and edges and what the graph model predicts.

\WhatToSay
This slide separates node features, edge attributes, and output fields. Nodes store coordinates, material parameters, scenario condition, and current dynamic state. Edges store relative spatial relations between neighbouring nodes. The output is the predicted next physical field vector for each mesh node. This makes MeshGraphNet suitable for state-transition style prediction. The key distinction from PINN is that physical quantities enter through features and training data, while the governing equations are not enforced as explicit residuals in the same way.

\KeyPoints
\begin{itemize}
  \item Nodes contain geometry, material data, and physical state.
  \item Edges describe local spatial relationships.
  \item The model predicts next physical fields over mesh nodes.
\end{itemize}

\Transition
Now I will summarize the graph neural network architecture.
\end{slidenote}

\begin{slidenote}{MeshGraphNet: Encoder--Processor--Decoder}
\Purpose
Describe the MeshGraphNet architecture without excessive implementation detail.

\WhatToSay
The architecture has three main parts. The encoder maps raw node and edge features into latent vectors. The processor applies repeated message-passing blocks over the graph connectivity. The decoder maps the final latent node states into predicted physical fields. In this implementation the model uses a latent dimension of $128$, ten message-passing steps, MLP encoders, and a node-wise decoder. This design is common in neural simulators because it separates feature embedding, learned interaction modelling, and output prediction.

\KeyPoints
\begin{itemize}
  \item Encoder: embeds raw node and edge features.
  \item Processor: performs repeated message passing.
  \item Decoder: converts latent node states to physical outputs.
\end{itemize}

\Transition
The next slide explains the message-passing block itself.
\end{slidenote}

\begin{slidenote}{MeshGraphNet: Message Passing Block}
\Purpose
Connect graph message passing with local physical interaction learning.

\WhatToSay
In the message-passing block, edge representations are updated using information from neighbouring node states and current edge attributes. Then incoming edge messages are accumulated for each destination node using scatter-add aggregation. Finally, the node state is updated with a residual connection and layer normalization. This can be interpreted as a learned local interaction mechanism on the mesh. It is useful for modelling spatial dependencies, but physical reliability still depends on validation against reference fields because the equations are learned implicitly rather than enforced as residuals.

\KeyPoints
\begin{itemize}
  \item Edge updates exchange information between neighbouring nodes.
  \item Aggregation collects incoming messages at each node.
  \item Physical behaviour is learned from data and must be validated.
\end{itemize}

\Transition
The fourth model family is the Transformer-style surrogate.
\end{slidenote}

\begin{slidenote}{Transformer: Inputs and Outputs}
\Purpose
Introduce the attention-based route and its token representation.

\WhatToSay
The Transformer route treats sampled points as tokens. Input tokens include coordinates, initial temperature and displacement state, velocity-like components, and material parameters. Query tokens contain the coordinates where fields are requested. The output at each query point is temperature and two in-plane displacement components. This token structure makes the route flexible because the number of input points and query points can differ. Its advantage is fast query-style prediction, while its limitation is weaker explicit physics enforcement compared with PINN.

\KeyPoints
\begin{itemize}
  \item Input tokens contain state, geometry, and material information.
  \item Query tokens specify prediction locations.
  \item Outputs are temperature and in-plane displacement at query points.
\end{itemize}

\Transition
Next, I will show how the Transformer architecture processes those tokens.
\end{slidenote}

\begin{slidenote}{Transformer: Architecture}
\Purpose
Summarize the encoder--decoder attention mechanism for query prediction.

\WhatToSay
The Transformer architecture has an encoder and a decoder. The encoder processes input tokens using self-attention layers, building a representation of the known thermoelastic state. The decoder processes query tokens and uses cross-attention to attend to the encoded input representation. Finally, a linear head returns the predicted temperature and displacement components. In intuitive terms, each requested query point can use attention to decide which input points are most relevant for its prediction. This supports flexible query-point inference, but attention is still a learned statistical mechanism rather than an explicit thermoelastic conservation law.

\KeyPoints
\begin{itemize}
  \item Encoder self-attention processes known input tokens.
  \item Decoder cross-attention connects query points with encoded inputs.
  \item The output head predicts $(T,u,v)$ at requested points.
\end{itemize}

\Transition
The next slide defines the training loss for this route.
\end{slidenote}

\begin{slidenote}{Transformer: Training Loss}
\Purpose
Explain the supervised objective and its limited physical regularisation.

\WhatToSay
The Transformer loss combines a supervised reconstruction term, a velocity consistency term, and a thermal residual term. The supervised loss reconstructs the COMSOL reference fields. The velocity consistency term aligns displacement rate with exported velocity. The thermal residual encourages diffusion-like behaviour through thermal diffusivity $k/(\rho C_p)$. The training uses AdamW and token subsampling with $N_{\mathrm{in}}=N_q=1024$ to control computational cost. The interpretation remains cautious: this is mainly a supervised attention-based surrogate, not a strict physics-informed solver like the PINN route.

\KeyPoints
\begin{itemize}
  \item The objective combines supervised, velocity, and thermal terms.
  \item Token subsampling controls attention cost.
  \item Physics constraints are weaker than in PINN.
\end{itemize}

\Transition
Now I will move from model descriptions to calculation status and results.
\end{slidenote}

\begin{slidenote}{Calculation Status}
\Purpose
Summarize the evaluation setup and state caveats before showing plots.

\WhatToSay
The calculation setup uses forty scenarios per material and four materials: granite, sandstone, limestone, and basalt. With four model routes, this gives $40\times4\times4=640$ model responses under one shared contract. The slide also emphasizes zero fallback: every plotted response is an active model output, not a mock result. The operating envelope remains two-dimensional plane strain. The FNO caveat is repeated here because it affects interpretation: FNO is included in the comparison, but its output scales are still miscalibrated, so it should currently be treated as a diagnostic branch.

\KeyPoints
\begin{itemize}
  \item The comparison covers $640$ model responses.
  \item Plotted outputs come from active models.
  \item FNO is included but currently miscalibrated.
\end{itemize}

\Transition
The first result slide compares predicted travel times.
\end{slidenote}

\begin{slidenote}{Travel-Time Predictions}
\Purpose
Explain what the travel-time plot shows and how to interpret the model differences.

\WhatToSay
This plot compares predicted travel time by model and material. The approximate values shown on the slide are about $0.11$ ms for PINN, $0.14$ ms for Transformer, $1.29$ ms for FNO, and $2.47$ ms for MeshGraphNet. The key observation is that model-family differences are larger than the basalt--sandstone material difference in the current prototype. This does not mean one model is universally correct. It means that model formulation strongly affects directional travel-time estimates, so comparison and validation are necessary before drawing physical conclusions.

\KeyPoints
\begin{itemize}
  \item The plot compares travel-time estimates across models and materials.
  \item Model-family differences dominate the material difference here.
  \item The result motivates careful validation.
\end{itemize}

\Transition
To evaluate accuracy, the next slide compares predictions with an analytical reference.
\end{slidenote}

\begin{slidenote}{Accuracy Against Analytical Reference}
\Purpose
Present the main quantitative result and its cautious interpretation.

\WhatToSay
Here the predicted travel times are compared with a simple analytical reference, $t_{\mathrm{gt}}=r/V_p$. In this formula, $t_{\mathrm{gt}}$ is reference travel time, $r$ is source--probe distance, and $V_p$ is compressional wave velocity. This is not a full thermoelastic validation; it is a sanity check under a simple isotropic approximation. In the balanced sweep, PINN has a median error of about $2.0\%$ and remains within the $10\%$ band in all cases. Transformer has a larger median error, while FNO and MeshGraphNet show much larger deviations in the current setup.

\KeyPoints
\begin{itemize}
  \item The analytical reference is a sanity check, not complete validation.
  \item PINN shows the best agreement with the reference.
  \item Other routes need calibration or stronger time conditioning.
\end{itemize}

\Transition
Accuracy is one side of the comparison; inference latency is also important.
\end{slidenote}

\begin{slidenote}{Inference Latency and Time Behavior}
\Purpose
Compare practical response time and explain the fixed-horizon issue.

\WhatToSay
The latency plot shows that all routes return predictions in milliseconds on the development machine, which is fast enough for interactive prototype use. The median latency is approximately $1.3$ ms for Transformer, $1.5$ ms for MeshGraphNet, $2.4$ ms for FNO, and $2.6$ ms for PINN. However, speed is not the only criterion. PINN is also the only route that changes smoothly with requested observation time. The other routes currently behave more like fixed-horizon predictors, so their temporal behaviour needs additional development for time-dependent directional studies.

\KeyPoints
\begin{itemize}
  \item All routes are fast enough for interactive prototype comparison.
  \item Transformer has the lowest median latency.
  \item PINN currently has the most meaningful time-conditioned behaviour.
\end{itemize}

\Transition
Before concluding, I will state the main threats to validity.
\end{slidenote}

\begin{slidenote}{Threats to Validity}
\Purpose
Show scientific caution and define the limits of the results.

\WhatToSay
The main threats to validity are listed here. First, the training data are limited and come from a pilot COMSOL-derived configuration, not from a full calibrated geological database. Second, the physical setup is simplified to two-dimensional plane strain with effective material parameters. Third, the current FNO route has a normalization and output-scale issue. Fourth, the analytical reference $r/V_p$ is useful but isotropic and incomplete for full thermoelastic validation. These limitations are important because the results should be interpreted only within the simplified assumptions of the study.

\KeyPoints
\begin{itemize}
  \item The dataset is limited and synthetic.
  \item The physical formulation is simplified.
  \item The analytical reference is useful but limited.
\end{itemize}

\Transition
With these limitations in mind, I can summarize the conclusions.
\end{slidenote}

\begin{slidenote}{Conclusions}
\Purpose
State the main scientific outcomes of the diploma work.

\WhatToSay
The thesis develops a physically grounded comparative framework for AI-guided prediction of thermoelastic wave propagation in geological media. Within the tested assumptions, PINN is the strongest physics-informed baseline and shows the best agreement with the analytical travel-time reference. Transformer is the fastest route and is useful for rapid query-style inference. MeshGraphNet remains a useful mesh-aware surrogate candidate. FNO should currently be treated as a diagnostic branch until output normalization is corrected. Overall, the work demonstrates a careful prototype for comparing neural surrogate approaches rather than a final geophysical simulator.

\KeyPoints
\begin{itemize}
  \item The main contribution is a comparative AI-assisted framework.
  \item PINN performs best against the travel-time sanity check.
  \item Other routes remain useful candidates with specific limitations.
\end{itemize}

\Transition
Finally, I will outline the most important next steps.
\end{slidenote}

\begin{slidenote}{Future Work}
\Purpose
Show how the prototype can be developed after the diploma.

\WhatToSay
The first future direction is to retrain and recalibrate the two-dimensional FNO route. The second is to add time-conditioned rollout behaviour to MeshGraphNet. The third is to extend the material catalog with more complete thermoelastic parameters and validation data. The fourth is to compare against additional analytical or numerical baselines. The final message is that the prototype estimates directional characteristics of thermoelastic wave propagation within simplified assumptions and provides a basis for further research.

\KeyPoints
\begin{itemize}
  \item Recalibrate FNO.
  \item Improve time-conditioned behaviour in MeshGraphNet.
  \item Extend material data and validation baselines.
\end{itemize}

\Transition
This completes the main presentation.
\end{slidenote}

\begin{slidenote}{Thank You}
\Purpose
Close the presentation and invite questions.

\WhatToSay
Thank you for your attention. I will be happy to answer your questions about the physical assumptions, the COMSOL dataset, the model comparison, or the interpretation of the results.

\KeyPoints
\begin{itemize}
  \item Close politely.
  \item Signal readiness to discuss assumptions and limitations.
  \item Keep answers consistent with the research-prototype framing.
\end{itemize}

\Transition
Questions from the committee.
\end{slidenote}

\begin{slidenote}{Backup: Mechanical Rock Parameters}
\Purpose
Use this backup only if the committee asks about mechanical material values.

\WhatToSay
This backup table gives the extended mechanical material catalogue. It includes density, Young's modulus, Poisson's ratio, shear modulus, bulk modulus, and representative P-wave and S-wave velocities for ten rock types. I would not read every value. The main point is that the prototype uses physically meaningful effective parameters, while still treating them as midpoint values under isotropic assumptions.

\KeyPoints
\begin{itemize}
  \item Mechanical parameters control inertia, stiffness, and wave speeds.
  \item Values are effective and comparative, not site-calibrated.
  \item The table supports the material assumptions used in the main flow.
\end{itemize}

\Transition
If needed, the next backup slide gives thermal and porosity parameters.
\end{slidenote}

\begin{slidenote}{Backup: Thermal and Porosity Parameters}
\Purpose
Use this backup if asked how thermal properties and porosity enter the catalogue.

\WhatToSay
This table gives thermal expansion, thermal conductivity, heat capacity, volumetric heat capacity, thermoelastic coupling, and porosity information. These quantities control how a temperature disturbance spreads and how temperature change contributes to thermal strain and stress. Dashes indicate values that were not available in the source catalogue. They are left as missing values rather than filled by unsupported assumptions.

\KeyPoints
\begin{itemize}
  \item Thermal properties control heat spreading and thermoelastic strain.
  \item Porosity is included as an effective descriptive property.
  \item Missing values are shown explicitly.
\end{itemize}

\Transition
If the question concerns wave velocities, use the next backup slide.
\end{slidenote}

\begin{slidenote}{Backup: Elastic Wave Velocities}
\Purpose
Explain the P-wave and S-wave formulas if the committee asks about wave physics.

\WhatToSay
This slide gives the standard formulas for compressional and shear wave velocities. P-waves involve particle motion mainly parallel to the propagation direction and include volume change. S-waves involve transverse motion and depend on shear rigidity. The formulas show that stiffness increases wave velocity, while density increases inertia. These relations also explain why the analytical reference $t=r/V_p$ is a useful first sanity check for travel time.

\KeyPoints
\begin{itemize}
  \item $V_P$ depends on bulk modulus, shear modulus, and density.
  \item $V_S$ depends on shear modulus and density.
  \item These formulas justify the simple travel-time reference.
\end{itemize}

\Transition
If the question concerns heat-spreading time scale, use the next backup slide.
\end{slidenote}

\begin{slidenote}{Backup: Thermal Diffusivity}
\Purpose
Explain the compact thermal parameter used to interpret heat spreading.

\WhatToSay
Thermal diffusivity is defined as $\alpha_T=k/(\rho C_p)$. It has units of square metres per second and describes how quickly a temperature disturbance spreads through a material. Larger thermal diffusivity means faster temperature smoothing. This quantity is useful because it combines conductivity, density, and heat capacity into one interpretable parameter for the thermal part of the coupled response.

\KeyPoints
\begin{itemize}
  \item Thermal diffusivity describes heat-spreading rate.
  \item It combines conductivity, density, and heat capacity.
  \item It helps interpret the time scale of the thermal response.
\end{itemize}

\Transition
The final backup slides document image credits and references.
\end{slidenote}

\begin{slidenote}{Image Credits}
\Purpose
Document visual sources if the committee asks about images.

\WhatToSay
This slide lists the sources for the external educational diagrams and rock textures used in the presentation. It helps separate illustrative visual material from the thesis results. The COMSOL field images and model-result plots are part of the project workflow, while the external diagrams are used only to support physical explanation.

\KeyPoints
\begin{itemize}
  \item External images are credited separately.
  \item Illustrative diagrams are not model outputs.
  \item Project-generated figures are distinct from source images.
\end{itemize}

\Transition
The reference slide gives the scientific and technical sources behind the work.
\end{slidenote}

\begin{slidenote}{References}
\Purpose
Point to the scientific basis of the thesis if asked.

\WhatToSay
The references include rock physics, thermoelasticity, geodynamics, surface geophysics, anisotropy, neural operators, and Transformer-related sources. In the defence, I would use this slide only in the question period. It supports the theoretical background, material-property assumptions, and machine-learning model choices used throughout the thesis.

\KeyPoints
\begin{itemize}
  \item References cover physics, geology, and neural surrogate modelling.
  \item They support the theory and material assumptions.
  \item The slide is mainly for Q\&A support.
\end{itemize}

\Transition
End of backup material.
\end{slidenote}

\newpage
\section*{\color{AcademicNavy}Possible questions from the committee}
\addcontentsline{toc}{section}{Possible questions from the committee}

\textbf{1. Why did you choose directional prediction instead of full-field simulation?}

Because the thesis focuses on a narrower and more defensible research task. Full-field thermoelastic simulation is more computationally and physically complex. Directional prediction allows us to compare source--probe responses, travel time, and model behaviour under controlled assumptions.

\textbf{2. Is the prototype validated on real field data?}

No. The current work uses COMSOL-derived synthetic data and analytical sanity checks. The results should be interpreted within simplified modelling assumptions. Field data would be an important next step, but it is outside the current scope.

\textbf{3. Why is PINN called a physics-informed baseline?}

PINN is called a physics-informed baseline because its training objective includes residual terms connected with thermoelastic relations. This means that physical structure is included in the learning objective, but the result is still interpreted as a baseline model inside the simplified prototype.

\textbf{4. Why does FNO perform poorly in the current results?}

The current two-dimensional FNO route has a calibration and output-scale issue. Therefore, it is kept in the comparison as a diagnostic branch. Its architecture is still relevant, but this implementation needs retraining and normalization correction.

\textbf{5. What is the role of the analytical reference $t=r/V_p$?}

It is a simple travel-time sanity check based on distance and P-wave velocity. It is not a complete thermoelastic validation, but it helps identify whether predicted travel times are physically plausible in a basic isotropic approximation.

\textbf{6. Why is Transformer fast but less physically reliable?}

The Transformer route uses attention-based query prediction and can return estimates quickly, but most of its behaviour is learned from data. It has weaker explicit physical constraints than PINN, so its predictions require careful validation.

\textbf{7. What is the current limitation of MeshGraphNet?}

MeshGraphNet preserves mesh topology and is suitable for mesh-aware surrogate modelling, but in the current prototype it behaves mainly as a fixed-horizon predictor. It needs stronger time-conditioned rollout behaviour for more reliable time-dependent prediction.

\textbf{8. What are the main limitations of the work?}

The main limitations are the simplified two-dimensional setting, effective isotropic material parameters, limited synthetic data, fixed-horizon behaviour in some surrogate routes, the FNO normalization issue, and the absence of direct calibration against real geological measurements.

\textbf{9. What is the main contribution of the diploma work?}

The main contribution is a physically grounded comparative framework for AI-Guided Prediction of Thermoelastic Wave Propagation in Geological Media. It combines thermoelastic theory, geological material parameters, a controlled dataset, and several neural surrogate routes under one consistent comparison workflow.

\newpage
\section*{\color{AcademicNavy}Short version of the speech}
\addcontentsline{toc}{section}{Short version of the speech}

Good morning. The topic of our diploma work is \textit{AI-Guided Prediction of Thermoelastic Wave Propagation in Geological Media}. The work presents a research prototype for AI-guided directional prediction of thermoelastic wave propagation in simplified geological media. It is not a field-calibrated simulator. The goal is to estimate selected directional characteristics and compare several neural surrogate approaches under one shared contract.

The physical motivation is that temperature changes in rocks can create thermal expansion, stress, displacement, and wave-like mechanical response. This is relevant for geothermal systems, underground engineering, and geophysical interpretation. A full coupled numerical simulation can be expensive when many materials and source--probe directions must be compared, so the thesis formulates a narrower directional prediction task.

The model uses effective geological material parameters such as density, elastic moduli, thermal conductivity, heat capacity, thermal expansion, and P-wave and S-wave velocities. The formulation is two-dimensional, assumes plane strain, small deformation, and linear isotropic thermoelasticity, and does not explicitly resolve pores, fluids, fractures, plasticity, or full anisotropy.

Directional prediction is defined by a source point and a probe point. Their difference gives the direction, distance, and azimuth. The prototype estimates temperature, displacement, travel time, and directional response for that query. The data are based on a controlled COMSOL-derived two-dimensional thermoelastic setup, exported and normalized for several model families.

The comparison includes a Physics-Informed Neural Network, Fourier Neural Operator, MeshGraphNet, and Transformer-style surrogate. PINN is used as the physics-informed baseline because its objective includes thermoelastic residual terms. FNO, MeshGraphNet, and Transformer-style models are treated as surrogate routes for comparison, not as physically equivalent numerical solvers.

The calculations use forty scenarios per material, four materials, and four model routes, giving six hundred forty model responses. In the travel-time comparison, model-family differences are larger than the material difference in the current prototype. Against the simple analytical reference $t=r/V_p$, PINN shows the best agreement and remains within the ten percent band in the balanced sweep. Transformer-style prediction is the fastest route, while MeshGraphNet and FNO need additional development for time behaviour and calibration.

The main limitations are the limited synthetic dataset, the simplified two-dimensional assumptions, the current FNO normalization issue, fixed-horizon behaviour in some routes, and the fact that the analytical reference is only a sanity check. The conclusion is that the work provides a physically grounded comparative framework for AI-guided thermoelastic wave prediction within simplified assumptions. Future work should recalibrate FNO, improve time-conditioned MeshGraphNet behaviour, extend material data, and compare against additional analytical or numerical baselines.

\end{document}
